\documentclass[11pt,a4paper]{article}
\usepackage[utf8]{inputenc}
\usepackage[T1]{fontenc}
\usepackage{amsmath,amssymb}
\usepackage[margin=1in]{geometry}
\usepackage{hyperref}

\title{IOI 2023 -- Closing Time}
\author{Solution}
\date{}

\begin{document}
\maketitle

\section{Problem Summary}
Given a weighted tree, two special vertices $X$ and $Y$, and budget $K$, choose closing times to maximize the total number of vertices reachable from $X$ and from $Y$.

\section{Solution}

\subsection{Two Costs per Vertex}
Let
\[
p_v = \min(\mathrm{dist}(X,v), \mathrm{dist}(Y,v)), \qquad
q_v = \max(\mathrm{dist}(X,v), \mathrm{dist}(Y,v)).
\]
Each vertex can contribute in two ways:
\begin{itemize}
  \item paying $p_v$ means it is covered from the nearer endpoint;
  \item paying $q_v$ means it is covered from both sides / after the common part is chosen.
\end{itemize}

\subsection{Case 1: Solutions with a Common Part}
If the chosen reachable sets share vertices on the path from $X$ to $Y$, the official solution sweeps a border along that path. Vertices already forced into the common part contribute fixed cost, while the remaining vertices are optimized using Fenwick trees over a sorted multiset of candidate values. This allows us to compute the best leftover score for each border in $O(\log n)$.

\subsection{Case 2: Solutions without a Common Part}
If the two reachable sets are disjoint, then each vertex contributes either $p_v$ or $q_v$, independently. The optimum is obtained by sorting all $2n$ candidate costs and greedily taking the cheapest ones within budget.

\subsection{Final Answer}
Take the maximum over the two cases above. Distances from $X$ and $Y$ are computed once with DFS, so the whole solution is dominated by sorting and Fenwick operations.

\section{Complexity}
\begin{itemize}
  \item Distance preprocessing: $O(n)$.
  \item Fenwick-based optimization and sorting: $O(n \log n)$.
  \item Total: $O(n \log n)$.
\end{itemize}

\end{document}
